\documentclass[11pt]{article}
\usepackage[margin=1in]{geometry}
\usepackage{amsmath,amssymb,mathtools}
\usepackage{fontspec,unicode-math}
\setmainfont{Latin Modern Roman}
\setsansfont{Latin Modern Sans}
\setmathfont{Latin Modern Math}
\usepackage{xurl,hyperref}
\hypersetup{colorlinks=true,urlcolor=blue,linkcolor=blue}
\newcommand{\sref}[2]{\hyperlink{source:#1:#2}{[S#2]}}
\newcommand{\cataloguescope}[1]{\paragraph{Recorded mathematical question.}#1}
\begin{document}
\section{Open Problems on Billiards and Geometric Optics}\label{statement}

\subsection{Definitions and mathematical statement}

\paragraph{Shared notation.}$\mathbb N=\{1,2,\ldots\}$, and $\mathbb Z,\mathbb Q,\mathbb R,\mathbb C$ denote the integers, rationals, reals and complex numbers. Any other conventions are specified locally.


The space of oriented lines has its standard symplectic structure, identified with the cotangent bundle of the sphere of directions. The open set $\Pi$ consists of lines crossing the strictly convex hypersurface twice transversely. A Hamiltonian symplectomorphism is the time-one transformation of a smooth Hamiltonian isotopy. $C^k$ proximity means closeness of hypersurfaces and their derivatives through order $k$; $C^\infty$ convergence means convergence of every derivative on compact subsets where maps are defined. The allowed compositions contain reflections only; inverses are not added as extra generators.
\paragraph{Statement recorded in the supplied source.}
Let $\gamma\subset\mathbb{R}^n$ be a closed strictly convex hypersurface, and let $\Pi$ be the phase cylinder of oriented lines meeting $\gamma$ transversely twice. For every $\varepsilon>0$ and $k\in\mathbb{N}$, is every $C^\infty$ Hamiltonian symplectomorphism $\Pi\to\Pi$ a $C^\infty$ limit of compositions of reflections from $\gamma$ and from hypersurfaces $\varepsilon$-close to $\gamma$ in the $C^k$ topology? \sref{statement}{2}
\subsection{Short English statement}
Can every smooth Hamiltonian map on billiard phase space be approximated smoothly by composing reflections in the original surface and arbitrarily small prescribed smooth perturbations, without adding inverse reflections?
\subsection{Sources}
\begin{description}
\item[\hypertarget{source:statement:1}{S1}] ULAM.AI and the UnsolvedMath Contributors, \emph{UnsolvedMath: A Curated Collection of Open Mathematics Problems}, catalogue entry. CC BY 4.0. \url{https://huggingface.co/datasets/ulamai/UnsolvedMath}.
\item[\hypertarget{source:statement:2}{S2}] M. Bialy, C. Fierobe, A. Glutsyuk, M. Levi, A. Plakhov and S. Tabachnikov, \emph{Open problems on billiards and geometric optics} (2021), arXiv:2110.10750v2, Alexey Glutsyuk, Problem 2 and remarks, pp. 6--7. \url{https://arxiv.org/pdf/2110.10750v2}.
\end{description}
\label{end:statement}
\end{document}
